% !TEX root = ./ECEN-601 - Homework 6 - Brody Jordan.tex

\documentclass[11pt]{article}

\usepackage[
    assignmentDueDate={November 18, 2025 at 11:59 pm}
]{C:/Users/bljor/Sync/Courses/assignment}

% make span a math command
\DeclareMathOperator{\spn}{span}

\begin{document}
\makeAssignmentTitle

% ------ PROBLEM 1 ------ %
\begin{homeworkProblem}

    Consider \R[4] with the standard inner product. Let W be the subspace of \R[4] consisting of all
    vectors which are orthogonal to both $\alpha = (1, 0, -1, 1)$ and $\beta = (2, 3, -1, 2)$. Find a basis for
    $W$. \\

    \solution
    We want to find the set of vectors $x = (x_1, x_2, x_3, x_4)$ such that $\langle x, \alpha \rangle = 0$ and $\langle x, \beta \rangle = 0$.
    This gives us the system of equations:
    \begin{align*}
        x_1 - x_3 + x_4 &= 0 \\
        2x_1 + 3x_2 - x_3 + 2x_4 &= 0
    \end{align*}
    From the first equation, we have $x_3 = x_1 + x_4$. Substituting this into the second equation:
    \begin{align*}
        2x_1 + 3x_2 - (x_1 + x_4) + 2x_4 &= 0 \\
        x_1 + 3x_2 + x_4 &= 0 \implies x_1 = -3x_2 - x_4
    \end{align*}
    Substituting back into the expression for $x_3$:
    \[
        x_3 = (-3x_2 - x_4) + x_4 = -3x_2
    \]
    So the general solution is:
    \[
        x = \begin{bmatrix} -3x_2 - x_4 \\ x_2 \\ -3x_2 \\ x_4 \end{bmatrix} = x_2 \begin{bmatrix} -3 \\ 1 \\ -3 \\ 0 \end{bmatrix} + x_4 \begin{bmatrix} -1 \\ 0 \\ 0 \\ 1 \end{bmatrix}
    \]
    Thus, a basis for $W$ is $\{(-3, 1, -3, 0), (-1, 0, 0, 1)\}$.

\end{homeworkProblem}

\pagebreak

% ------ PROBLEM 2 ------ %
\begin{homeworkProblem}

    Let $V$ be an inner product space, and let $\alpha$ and $\beta$ be vectors in $V$. Show that $\alpha = \beta$ if and
    only if $\langle\alpha \mid \gamma\rangle = \langle\beta \mid \gamma\rangle$ for every $\gamma$ in $V$. \\

    \solution

\end{homeworkProblem}

\pagebreak

% ------ PROBLEM 3 ------ %
\begin{homeworkProblem}

    Let
    \begin{equation*}
        \mathbf{p}_1 =
        \begin{bmatrix}
            1 \\
            2 \\
            3 \\
            4 \\
        \end{bmatrix} \quad
        \mathbf{p}_2 =
        \begin{bmatrix}
            4  \\
            -2 \\
            -6 \\
            -7 \\
        \end{bmatrix} \quad
        \mathbf{p}_3 =
        \begin{bmatrix}
            3  \\
            4  \\
            -2 \\
            1  \\
        \end{bmatrix}
    \end{equation*}
    and
    \[
        \mathbf{x} = \begin{bmatrix}
            1 \\
            2 \\
            3 \\
            7 \\
        \end{bmatrix}.
    \]
    Determine the best approximation, $\hat{\mathbf{x}}$, of the vector $\mathbf{x}$ by vectors in $W = \spn\{\mathbf{p}_1, \mathbf{p}_2, \mathbf{p}_3\}$.
    Also, determine the projection of $\mathbf{x}$ onto the orthogonal complement of $W$. \\

    \solution


\end{homeworkProblem}

\pagebreak

% ------ PROBLEM 4 ------ %
\begin{homeworkProblem}

    Let $S$ be a subset of an inner product space $V$. Show that $(S^\perp)^\perp$ contains the subspace
    spanned by $S$. When $V$ is finite-dimensional, show that $(S^\perp)^\perp$ is the subspace spanned by $S$. \\

    \solution

\end{homeworkProblem}

\pagebreak

% ------ PROBLEM 5 ------ %
\begin{homeworkProblem}

    Let $V$ be a finite-dimensional inner product space, and let $\mathcal{B} = (\alpha_1, \ldots, \alpha_n)$ be an orthonormal
    basis for $V$. Let $T$ be a linear operator on $V$; and $A$, the matrix of $T$ in the ordered basis $\mathcal{B}$.
    Prove that
    \[
        A_{ij} = \langle T \alpha_j \mid \alpha_i \rangle
    \]

    \solution

\end{homeworkProblem}

\pagebreak

% ------ PROBLEM 6 ------ %
\begin{homeworkProblem}

    Suppose $V = W_1 \oplus W_2$ and that $f_1$ and $f_2$ are inner products on $W_1$ and $W_2$, respectively.
    Show that there is a unique inner product $f$ on $V$ such that
    \begin{enumerate}[label=(\alph*), topsep=5pt, itemsep=5pt]
        \item $W_2 = W_1^\perp$;
        \item $f(\alpha, \beta) = f_k(\alpha, \beta)$, when $\alpha$, $\beta$ are in $W_k, k = 1, 2$. \newline
    \end{enumerate}

    \solution

\end{homeworkProblem}


\end{document}